% TEMPLATE-MILC-v3.tex - plantilla maestra UNICA de las guias MILC v3.
%
% La usan las cuatro familias de guias, cada una con su driver:
%   · Grados 6-11          -> scripts/build-guias-g11.py
%   · Bebras               -> scripts/build-guias-bebras.py
%   · Semillero            -> scripts/build-guias-semillero.py
%   · Territorio interior  -> scripts/build-guias-territorio-interior.py
%
% Antes habia tres copias casi identicas (~400 lineas iguales, 6 distintas):
% cada mejora habia que aplicarla tres veces, y en la practica no se hacia
% (el motor de recursos vivia solo en dos de ellas). Lo que varia entre
% familias esta parametrizado en 6 placeholders que cada driver llena
% (se nombran aqui SIN su sintaxis real: el builder reemplaza por texto plano,
% tambien dentro de los comentarios, y un valor multilinea rompe el preambulo):
%
%   HEADER_IZQ        encabezado izquierdo de pagina
%   HEADER_DER        encabezado derecho de pagina
%   PORTADA_ETIQUETA  rotulo sobre el numero grande (GRADO/MOMENTO/MODULO)
%   PORTADA_SERIE     linea de serie de la portada
%   PORTADA_ID        identificador de la guia en la portada
%   PORTADA_CIRCULO   el circulito de grado (°), o vacio si no aplica
%   PDF_TITULO        titulo en texto plano (sin \textbf ni comillas TeX) para
%                     los metadatos del PDF (/Title); lo produce lib_xelatex.texto_pdf
%
% Composicion (auditoria 2026-09-03): las bandas de estacion piden espacio
% (\Needspace*) y prohiben el salto justo despues (\nopagebreak), asi no quedan
% huerfanas al pie; las cajas largas (softbox, drawbox, infoband, puentebox) son
% `breakable`, asi una caja de media pagina ya no arrastra todo a la siguiente
% dejando la anterior al 40 %. Todo texto sobre mostaza o verde va en milcNegro
% (blanco sobre #E5B400 da 1,93:1 y sobre #58A923 2,95:1; ninguno pasa WCAG AA).
% El resto de claves las documenta content/guias/_SCHEMA.md. El script valida
% que no queden placeholders sin reemplazar antes de escribir el archivo final.

% Etiquetado PDF (latex-lab): /Lang, estructura y texto alternativo de las
% figuras, para lectores de pantalla. Debe ir ANTES de \documentclass.
\DocumentMetadata{lang=es-CO,tagging=on}
\documentclass[11pt,letterpaper]{article}
% headheight=24pt: la cabecera derecha lleva dos lineas (identidad de la guia
% y estacion actual). headsep se acorta en la misma medida para que el bloque
% de texto quede exactamente donde estaba con headheight=12pt.
\usepackage[margin=2.0cm,headheight=24pt,headsep=13pt]{geometry}
\usepackage[table]{xcolor}
\usepackage{fontspec}
\usepackage[spanish]{babel}
\usepackage{tikz}
% Librerías TikZ para los diagramas de las guías (bloque `recursos:`):
%  · babel  → babel en español vuelve activos caracteres como «>», y TikZ los usa
%             en sus opciones (>=stealth, ->). Sin esta librería, cualquier
%             diagrama con flechas revienta con «\language@active@arg> has an
%             extra }». Debe ir ANTES de que se dibuje nada.
%  · positioning        → sintaxis «below left=2pt and 3pt».
%  · decorations.pathreplacing → llaves ({brace}) para señalar rangos.
\usetikzlibrary{babel,positioning,decorations.pathreplacing}
\usepackage{graphicx}
% Circuitos: el trabajo de electrónica se hace hoy con micro:bit y sus sensores
% integrados (MakeCode, sin cableado), así que no se necesitan esquemáticos.
% Si algún día entran montajes con protoboard, basta descomentar esta línea:
% los diagramas .tex/.tikz declarados en `recursos:` ya se incrustan solos.
% \usepackage{circuitikz}
\usepackage[most]{tcolorbox}
\usepackage{tabularx}
\usepackage{array}
\usepackage{fancyhdr}
\usepackage{titlesec}
\usepackage[document]{ragged2e}  % bandera a la derecha en todo el cuerpo
\usepackage{enumitem}
\usepackage{microtype}
\usepackage{etoolbox}
\usepackage{needspace}
% hyperref al final del preambulo: metadatos del PDF (titulo, autor, idioma,
% familia), marcadores por estacion (\pdfbookmark) y enlaces sin recuadro.
\usepackage[hidelinks,
  pdftitle={Sentir no es actuar: el espacio entre lo que siento y lo que hago},
  pdfauthor={PhD. Álvaro Cárdenas Orozco},
  pdfsubject={Territorio interior · Educación socioemocional · Grado 6° · Momento 3 · MILC v3},
  pdflang={es-CO},
  bookmarksopen=true,bookmarksnumbered=false]{hyperref}

% Helvetica Neue no trae la flecha U+2192, asi que cada «→» escrito literalmente
% en el YAML se compilaba a NADA, en silencio (462 flechas en 61 guias: rutas del
% tipo «paleta → tipografia → lockup» salian rotas). Mapeamos el caracter a su
% version matematica, que si tiene glifo. Asi el autor puede seguir escribiendo
% «→» comodamente en el YAML.
\usepackage{newunicodechar}
\newunicodechar{→}{\ensuremath{\rightarrow}}
% Mismo problema con los simbolos de marca y vinetas: el «✓» de «una palomita ✓»
% salia como cuadrito vacio en el PDF de todas las guias que lo usaban.
\usepackage{pifont}
\newunicodechar{✓}{\ding{51}}
\newunicodechar{✗}{\ding{55}}
\newunicodechar{✔}{\ding{52}}
\newunicodechar{•}{\textbullet}
\newunicodechar{≥}{\ensuremath{\geq}}
\newunicodechar{≤}{\ensuremath{\leq}}

\setmainfont{Helvetica Neue}
\setsansfont{Helvetica Neue}
\setlength{\parindent}{0pt}
\setlength{\parskip}{6pt}
% Sin particion de palabras: en 6.º-7.º una palabra cortada por guion al final
% de linea («Comuni-cacion») frena la lectura mucho mas que una linea corta.
\hyphenpenalty=10000
\exhyphenpenalty=10000
\pretolerance=2000
\tolerance=3000
\setlength{\emergencystretch}{3em}
\linespread{1.10}
\renewcommand{\arraystretch}{1.25}
\setlist{leftmargin=5mm,itemsep=1mm,topsep=1mm}
\newcolumntype{Y}{>{\RaggedRight\arraybackslash}X}

\definecolor{milcMagenta}{HTML}{E6007E}
\definecolor{milcVino}{HTML}{5A0038}
\definecolor{milcTurquesa}{HTML}{0093A5}
\definecolor{milcVerde}{HTML}{58A923}
\definecolor{milcMostaza}{HTML}{E5B400}
\definecolor{milcOcre}{HTML}{B86600}
\definecolor{milcNegro}{HTML}{241816}
\definecolor{milcGris}{HTML}{F7F7F4}
\definecolor{milcLinea}{HTML}{E8E2DD}
\definecolor{milcGradePrimary}{HTML}{0D9488}
\definecolor{milcGradeDark}{HTML}{115E59}
\definecolor{milcGradeSoft}{HTML}{CCFBF1}
\definecolor{milcGradeAccent}{HTML}{CFE7FF}
\definecolor{milcGradeText}{HTML}{FFFFFF}
\color{milcNegro}

\pagestyle{fancy}
\fancyhf{}
% Cabecera de dos lineas: identidad de la guia arriba y, a la derecha y en
% gris, la estacion en curso (\markright desde \milcestacion). Antes de la
% primera estacion la segunda linea va vacia; el \strut mantiene la altura
% para que la linea de arriba no baile entre paginas.
\lhead{\small Territorio interior · Educación socioemocional\\[-1pt]{\footnotesize\strut}}
\rhead{\small Grado 6° · Momento 3 · MILC v3\\[-1pt]{\footnotesize\strut\color{milcNegro!65}\rightmark}}
\cfoot{\small\thepage}
\renewcommand{\headrulewidth}{0pt}

% Sobre mostaza (#E5B400) y verde (#58A923) el texto va en milcNegro; sobre el
% resto de la paleta, en blanco. Se usa dentro de fonttitle/fontupper, donde un
% \color posterior manda sobre coltitle.
\newcommand{\milctintasobre}[1]{%
  \ifboolexpr{test{\ifstrequal{#1}{milcMostaza}} or test{\ifstrequal{#1}{milcVerde}}}%
    {\color{milcNegro}}{\color{white}}}

\newtcolorbox{titlebox}[1]{
  enhanced,colback=#1,colframe=#1,arc=7mm,boxrule=0pt,
  left=7mm,right=7mm,top=3mm,bottom=3mm,
  before skip=8pt,after skip=8pt,
  fontupper=\sffamily\bfseries\Large\color{white}
}

\newtcolorbox{softbox}[3]{
  enhanced,breakable,pad at break=3mm,
  colback=#2,colframe=#1,borderline west={4pt}{0pt}{#1},
  arc=2mm,boxrule=.4pt,left=4mm,right=4mm,top=3mm,bottom=3mm,
  before skip=6pt,after skip=8pt,title={#3},coltitle=white,
  colbacktitle=#1,fonttitle=\sffamily\bfseries\milctintasobre{#1},fontupper=\RaggedRight,
  attach boxed title to top left={xshift=3mm,yshift=-2mm},
  boxed title style={arc=2mm, boxrule=0pt}
}

\newtcolorbox{drawbox}[2]{
  enhanced,breakable,pad at break=4mm,
  colback=white,colframe=#1,arc=4mm,boxrule=1pt,
  left=5mm,right=5mm,top=4mm,bottom=4mm,before skip=8pt,after skip=8pt,
  title={#2},coltitle=white,colbacktitle=#1,fonttitle=\sffamily\bfseries\milctintasobre{#1},
  fontupper=\RaggedRight,
  attach boxed title to top left={xshift=4mm,yshift=-2mm},
  boxed title style={arc=3mm, boxrule=0pt}
}

\newtcolorbox{infoband}[1]{
  enhanced,breakable,pad at break=2mm,colback=#1!8,colframe=#1!50,arc=2mm,boxrule=.5pt,
  left=4mm,right=4mm,top=2mm,bottom=2mm,before skip=4pt,after skip=4pt,
  fontupper=\small\RaggedRight
}

% Puente narrativo entre secciones (contrato editorial Conéctate).
% Da continuidad: "Acabas de X. Ahora vas a Y porque Z."
% Visualmente: banda izquierda en color del grado, texto en itálica pequeña.
\newtcolorbox{puentebox}{
  enhanced,breakable,pad at break=2mm,colback=milcGradeSoft,colframe=milcGradeDark,
  borderline west={3pt}{0pt}{milcGradeDark},
  arc=0mm,boxrule=0pt,left=5mm,right=4mm,top=2.5mm,bottom=2.5mm,
  before skip=4pt,after skip=8pt,
  fontupper=\itshape\small\color{milcNegro}\RaggedRight
}
% La flecha va en modo matematico a proposito: Helvetica Neue NO tiene el glifo
% U+2192, asi que \textrightarrow se compilaba a NADA («Missing character: There
% is no → in font Helvetica Neue») y el puente salia sin flecha. En modo
% matematico la flecha viene de la fuente matematica y si se dibuja.
\newcommand{\puente}[1]{%
  \begin{puentebox}%
    \textcolor{milcGradeDark}{$\rightarrow$}\hspace{2mm}#1%
  \end{puentebox}%
}

\newcommand{\linead}[1]{\noindent\color{milcLinea}\rule{#1}{0.5pt}\color{milcNegro}}
\newcommand{\respuesta}[1]{\par\vspace{2mm}\foreach \n in {1,...,#1}{\noindent\color{milcLinea}\rule{\linewidth}{0.5pt}\par\vspace{5mm}}\color{milcNegro}}
\newcommand{\checkbox}{\textcolor{milcGradeDark}{\fbox{\phantom{x}}}\hspace{5pt}}

% ─── Listas aireadas para las guías socioemocionales ────────────────────
% Componer las verificaciones como «\checkbox texto\\» deja el texto que salta
% de línea debajo del cuadrito y apelmaza el bloque. Estos entornos alinean la
% continuación y airean los ítems. Los usa Territorio Interior; cualquier
% familia puede hacerlo.
\newlist{checklista}{itemize}{1}
\setlist[checklista]{
  label=\textcolor{milcGradeDark}{\fbox{\phantom{x}}},
  leftmargin=9mm, labelsep=3.2mm, itemsep=2.4mm,
  topsep=2.5mm, parsep=0pt, partopsep=0pt, before=\RaggedRight
}
\newlist{pasoslista}{enumerate}{1}
\setlist[pasoslista]{
  label=\textcolor{milcGradeDark}{\sffamily\bfseries\arabic*.},
  leftmargin=8mm, labelsep=3mm, itemsep=2.8mm,
  topsep=1mm, parsep=0pt, partopsep=0pt, before=\RaggedRight
}

\newcommand{\phasecard}[4]{
\begin{tcolorbox}[enhanced,colback=#3,colframe=#2,arc=3mm,boxrule=.5pt,height=3.05cm,valign=center,halign=center,left=1.5mm,right=1.5mm,top=2mm,bottom=2mm]
{\sffamily\bfseries\fontsize{22}{24}\selectfont\textcolor{#2}{#1}}\par
{\sffamily\bfseries\small #4}
\end{tcolorbox}
}

% ── Piezas del rediseno 2026 (legibilidad 6.º-11.º) ───────────────────
% Banda de estacion: reemplaza el titlebox macizo. Titulo a la izquierda,
% dato practico (tiempo/modalidad) a la derecha, en una sola linea.
% Nunca huerfana: pide 9 lineas libres (o salta de pagina rellenando con
% \vfill) y prohibe el salto justo despues de la banda. Nueve y no seis:
% la caja partible que sigue necesita ~80pt (titulo adosado, relleno y dos
% lineas) para arrancar en la misma pagina; con menos, tcolorbox la manda
% entera a la siguiente y la banda se queda sola. El penalty 10000 va
% en `after`, ANTES del espacio posterior: un `after skip` (aunque sea 0pt)
% es un \vskip tras la caja, y ese glue es un punto de corte legal que un
% \nopagebreak escrito despues de \end{tcolorbox} ya no cierra. Cada banda
% deja marcador PDF y actualiza la cabecera derecha.
\newcounter{milcestacion}
\newcommand{\milcestacion}[2]{%
\Needspace*{9\baselineskip}%
\stepcounter{milcestacion}%
\pdfbookmark[1]{#2}{milcestacion\themilcestacion}%
\markright{#2}%
\begin{tcolorbox}[enhanced,colback=#1,colframe=#1,arc=2mm,boxrule=0pt,
  left=5mm,right=5mm,top=2.6mm,bottom=2.6mm,before skip=11pt,
  after={\par\nopagebreak\vskip7pt}]
{\sffamily\bfseries\fontsize{15}{18}\selectfont\milctintasobre{#1}#2}
\end{tcolorbox}}

% Tarjeta de paso numerada: sustituye las tablas «Pilar / Aplicacion», que
% eran formato de consulta docente, no de estudiante. El numero va en un
% circulo sobre el borde para que la secuencia se lea de un vistazo.
\newcommand{\milcpaso}[3]{%
\Needspace{4\baselineskip}%
\begin{tcolorbox}[enhanced,colback=white,colframe=milcLinea,arc=1.5mm,boxrule=.9pt,
  left=14mm,right=4mm,top=3mm,bottom=3mm,before skip=4pt,after skip=4pt,
  fontupper=\RaggedRight,
  overlay={\node[anchor=north west,circle,fill=#2,text=white,inner sep=0pt,
    minimum size=7.4mm,font=\sffamily\bfseries\small]
    at ([xshift=3.6mm,yshift=-3.2mm]frame.north west) {#1};}]
#3
\end{tcolorbox}}

% Casilla marcable de tamano util con lapiz (la anterior era \fbox{\phantom{x}},
% de alto variable segun el texto que la seguia).
\newcommand{\milcmarca}{\framebox[3.8mm]{\rule{0pt}{3.4mm}}}

% Fila marcable de lista suelta (checks de la estacion 1).
\newcommand{\milccheckrow}[1]{%
\par\vspace{1.8mm}\noindent
\makebox[6.4mm][l]{\raisebox{-.55mm}{\textcolor{milcNegro}{\milcmarca}}}%
\parbox[t]{\dimexpr\linewidth-6.4mm\relax}{\RaggedRight #1}\par}

% Celda marcable dentro de la rubrica (tabla de autoevaluacion). Misma
% sangria francesa, pero pensada para vivir dentro de una columna X.
\newcommand{\milcopcion}[1]{%
\makebox[5.2mm][l]{\raisebox{-.55mm}{\textcolor{milcNegro}{\milcmarca}}}%
\parbox[t]{\dimexpr\linewidth-5.2mm\relax}{\RaggedRight #1}}

% ── Recursos visuales de la guía ──────────────────────────────────────
% Declarados en el bloque `recursos:` del YAML y emitidos por el builder.
% \guiaFigura   → imagen rasterizada o PDF (foto, carta celeste, montaje).
% \guiaDiagrama → fuente TikZ/circuitikz (.tex) incrustada como vector:
%                 es la vía para circuitos y esquemáticos editables.
% [#1] = texto alternativo (alt) · #2 = ruta relativa al .tex · #3 = pie
% El alt llega al PDF etiquetado (/Alt) y lo lee el lector de pantalla.
\newcommand{\guiaFigura}[3][]{%
\begin{center}
\includegraphics[width=0.88\linewidth,height=0.38\textheight,keepaspectratio,alt={#1}]{#2}\\[1.5mm]
{\footnotesize\itshape\textcolor{milcNegro!75}{#3}}
\end{center}
}

\newcommand{\guiaDiagrama}[2]{%
\begin{center}
\input{#1}\\[1.5mm]
{\footnotesize\itshape\textcolor{milcNegro!75}{#2}}
\end{center}
}

\begin{document}
\thispagestyle{empty}

% ── Portada ───────────────────────────────────────────────────────────
% Antes: un rectangulo de color a pagina completa. Se veia bien en pantalla,
% pero en la fotocopiadora del colegio se comia el toner y salia gris sucio.
% Ahora la banda ocupa el tercio superior y el resto va en blanco.
\begin{tikzpicture}[remember picture,overlay]
  \path[fill=milcGradePrimary]
    (current page.north west) rectangle ([yshift=-10.6cm]current page.north east);

  \node[anchor=north west,text=milcGradeText] at ([xshift=2.0cm,yshift=-1.55cm]current page.north west)
    {\sffamily\bfseries\fontsize{12}{14}\selectfont MOMENTO};
  \node[anchor=north west,text=milcGradeText,opacity=.85] at ([xshift=2.0cm,yshift=-2.35cm]current page.north west)
    {\sffamily\bfseries\fontsize{8.5}{10}\selectfont TERRITORIO INTERIOR · SOCIOEMOCIONAL};
  \node[anchor=north east,text=milcGradeText,opacity=.85,align=right] at ([xshift=-2.0cm,yshift=-1.55cm]current page.north east)
    {\sffamily\bfseries\fontsize{9}{12}\selectfont I.E. Sor María Juliana\\Cartago, Valle del Cauca};

  \node[anchor=west,text=milcGradeText] (gradonum) at ([xshift=1.85cm,yshift=-7.1cm]current page.north west)
    {\sffamily\bfseries\fontsize{112}{112}\selectfont 3};
  % El circulito de grado (°) solo aplica a las guias de grado («11°»); en
  % Bebras y Semillero el numero grande es un momento/modulo, no un grado.
  % Se posiciona relativo al nodo `gradonum`, no en coordenadas absolutas.
  

  \node[anchor=west,text=milcGradeText,text width=11.4cm,align=left] at ([xshift=1.5cm,yshift=0.5cm]gradonum.east)
    {
     {\sffamily\bfseries\fontsize{9}{11}\selectfont\MakeUppercase{Territorio interior · Socioemocional}}\\[3mm]
     {\sffamily\bfseries\fontsize{21}{25}\selectfont\setlength{\baselineskip}{27pt}Sentir no es actuar\\[4pt]el espacio entre\\[4pt]sentir y hacer}\\[3.5mm]
     {\sffamily\bfseries\fontsize{15}{18}\selectfont Grado 6° · Momento 3}
    };
\end{tikzpicture}

\vspace*{9.4cm}
\vfill

{\sffamily\bfseries\fontsize{8}{10}\selectfont\textcolor{milcGradeDark}{LO QUE TE LLEVAS HOY}}

\begin{tcolorbox}[enhanced,colback=white,colframe=milcNegro,arc=2mm,boxrule=1.6pt,
  left=5mm,right=5mm,top=4mm,bottom=4mm,before skip=3pt,after skip=12pt,
  fontupper=\RaggedRight\fontsize{11}{15.5}\selectfont]
Tu fogón: un procedimiento propio de tres pasos para los momentos en que la emoción quiere decidir por ti — la señal que te avisa, la pausa concreta que haces y la frase que dices antes de responder.
\end{tcolorbox}

\begin{tcolorbox}[enhanced,colback=milcGris,colframe=milcGris,arc=2mm,boxrule=0pt,
  left=5mm,right=5mm,top=3.5mm,bottom=3.5mm,after skip=12pt]
{\sffamily\bfseries\fontsize{8}{10}\selectfont\textcolor{milcNegro!55}{ESTA GUÍA ES DE}}\\[3mm]
\begin{tabularx}{\linewidth}{X p{3.2cm} p{3.2cm}}
\linead{\linewidth} & \linead{3.0cm} & \linead{3.0cm}\\[-1mm]
{\fontsize{8.5}{10}\selectfont\textcolor{milcNegro!55}{Nombre completo}} &
{\fontsize{8.5}{10}\selectfont\textcolor{milcNegro!55}{Grupo}} &
{\fontsize{8.5}{10}\selectfont\textcolor{milcNegro!55}{Fecha}}\\
\end{tabularx}
\end{tcolorbox}

\vfill

\noindent\textcolor{milcLinea}{\rule{\linewidth}{1pt}}\\[2mm]
{\sffamily\bfseries\fontsize{9.5}{12}\selectfont Autor: Dr. Álvaro Cárdenas Orozco}\\[1mm]
{\sffamily\fontsize{8.5}{11}\selectfont\textcolor{milcNegro!55}{Metodología MILC · Apertura ancestral · Regulación emocional · Triángulo Dussel-Estoicismo-Floridi}}

\newpage

\section*{Apertura: Sentir no es actuar: el espacio entre lo que siento y lo que hago}

\begin{softbox}{milcGradeDark}{milcGradeSoft}{Saber ancestral}
Entre el pueblo \textbf{misak} ---en Guambía, Cauca--- la casa tiene un centro, y ese centro es la cocina: el \textbf{nak chak}. Y la cocina tiene el suyo: el \textbf{nak kuk}, el fogón. Alrededor de ese fuego se sienta la familia, se enseña y se aprende. Los mayores lo dicen de una manera que vale la pena oír despacio: \emph{«el derecho nace de las cocinas»}, porque de ahí nace y se difunde el \textbf{consejo} ---en su lengua, el \textbf{kørøsrøp}---. \textbf{Fíjate en lo que eso significa en la práctica.} Cuando pasa algo difícil, el asunto \textbf{no se resuelve donde pasó}: se lleva al fogón. Ahí se cuenta, se escucha a los mayores, se le da vueltas mientras el fuego arde. La decisión llega \textbf{después}, y llega distinta. \emph{Ese rodeo no es lentitud: es el lugar donde el consejo alcanza a nacer.} \textbf{El fogón es, literalmente, un espacio puesto entre lo que ocurre y lo que se hace.} Eso mismo es lo que hoy vas a construir por dentro: tu propio espacio entre lo que sientes y lo que haces.
\end{softbox}

\begin{softbox}{milcTurquesa}{milcGris}{Saber contemporáneo}
Hay una frase que se repite mucho y que hace daño: \emph{«me hizo enojar y por eso le pegué»}. Ahí van tres cosas distintas metidas en una sola: la \textbf{emoción} (la rabia), el \textbf{impulso} (las ganas de pegar) y la \textbf{conducta} (pegar). \textbf{La emoción llega sola, el impulso también, pero la conducta se elige}: entre el impulso y la mano hay un espacio, aunque a los once años se sienta de un milímetro. Se puede ensanchar, y no con fuerza de voluntad. Con \textbf{palabras}: en 2007, un equipo dirigido por Matthew Lieberman midió con resonancia magnética qué pasa en el cerebro cuando alguien le pone nombre a lo que siente, y encontró que \textbf{ponerlo en palabras baja la actividad de la amígdala} ---la parte que dispara la alarma--- y sube la de la corteza prefrontal, la que decide (Lieberman y otros, 2007). \emph{Es decir: el muestrario que armaste en el momento anterior no era una tarea de vocabulario.} Era el freno. Nombrar lo que sientes, mientras lo sientes, es lo que abre el espacio.
\end{softbox}

\begin{softbox}{milcMostaza}{milcGris}{Pregunta-puente}
\textit{Los misak no resuelven en el patio lo que pasó en el patio: lo llevan al fogón, y allí nace el consejo. \textbf{¿Dónde está tu fogón}\ldots ese lugar o ese momento al que puedes llevar lo que sientes antes de decidir qué hacer con ello?}
\end{softbox}

\begin{softbox}{milcVerde}{milcGris}{Saber y saber hacer}
Al terminar podrás: (1) \textbf{identificar}, en situaciones tuyas de verdad, la diferencia entre lo que sentiste, lo que te dieron ganas de hacer y lo que hiciste; (2) \textbf{explicar} por qué una emoción informa pero no obliga, y qué le pasa a una emoción con el tiempo si no la alimentas; (3) \textbf{crear} tu procedimiento de tres pasos ---señal, pausa y frase--- para los momentos en que la emoción quiere decidir por ti; (4) \textbf{aplicarlo} en un ensayo con un compañero, antes de necesitarlo de verdad.
\end{softbox}



\small
\begin{tabularx}{\textwidth}{>{\bfseries}p{3.0cm}Y}
\rowcolor{milcGradePrimary!12}Autor & Dr. Álvaro Cárdenas Orozco\\
DBA & Comprende los fenómenos que afectan a su territorio, regula sus emociones ante ellos y acompaña a otros con empatía (transversal 6°--11°, articulado con la política «Escuela, territorio de vida» del MEN, Resolución 006519 de 2025, y la Ley 1620 de 2013)\\
Referentes & Primera ayuda psicológica (OMS, 2011/2012) · Directrices IASC (2007) · USGS y Servicio Geológico Colombiano · Pueblo nasa y la minga (CRIC) · Dussel · Estoicismo · Floridi\\
Producto final & Tu fogón: un procedimiento propio de tres pasos para los momentos en que la emoción quiere decidir por ti — la señal que te avisa, la pausa concreta que haces y la frase que dices antes de responder.\\
\end{tabularx}
\normalsize

\puente{En el momento anterior aprendiste a nombrar lo que sientes con precisión. Hoy vas a usar esas palabras para algo concreto: \textbf{frenar}. Mira la ruta: primero tus propios casos, después cómo funciona el espacio, y al final tu fogón.}

\milcestacion{milcGradeDark}{Ruta MILC: así vamos a aprender}
\begin{tabularx}{\textwidth}{>{\centering\arraybackslash}X>{\centering\arraybackslash}X>{\centering\arraybackslash}X>{\centering\arraybackslash}X}
\phasecard{1}{milcVerde}{milcGris}{Escucha\\[-1mm]\small Observo, escucho y pregunto.} &
\phasecard{2}{milcTurquesa}{milcGris}{Sistematización\\[-1mm]\small Distingo y comprendo.} &
\phasecard{3}{milcMagenta}{milcGris}{Praxis\\[-1mm]\small Construyo y aplico.} &
\phasecard{4}{milcVino}{milcGris}{Evaluación\\[-1mm]\small Reflexión liberadora.}
\end{tabularx}


\puente{Antes de entender nada, mira tus últimos días. No para juzgarte: para ver de qué tamaño es hoy tu espacio entre sentir y hacer.}

\milcestacion{milcVerde}{Estación 1 de 3 · Escucha}

\begin{softbox}{milcVerde}{milcGris}{Escena para escuchar}
\textbf{Actividad 1 · IDENTIFICA --- Tres casillas, tres veces.} \textbf{Nadie va a leer esta página.} Piensa en \textbf{tres momentos} de las últimas dos semanas en que reaccionaste de una manera que después no te gustó: un grito, un mensaje enviado con rabia, una salida de la mesa dando un portazo, un comentario que hirió, un silencio con el que castigaste a alguien. Para cada uno, tres casillas: \textbf{qué sentí} (usa tu muestrario del momento anterior, la palabra exacta), \textbf{qué me dieron ganas de hacer} y \textbf{qué hice}. Al final, una línea más: \emph{¿cuánto tiempo pasó entre lo segundo y lo tercero?} Segundos, minutos, nada. \emph{Ese número es el tamaño actual de tu espacio, y es lo que vamos a agrandar.}
\end{softbox}

{\sffamily\bfseries\fontsize{8}{10}\selectfont\textcolor{milcNegro!55}{MARCA CUANDO LO HAYAS HECHO}}
\milccheckrow{Escribiste tres situaciones reales de las últimas dos semanas, con la palabra exacta de lo que sentiste.}
\milccheckrow{Separaste en cada una lo que te dieron ganas de hacer de lo que efectivamente hiciste.}
\milccheckrow{Calculaste cuánto tiempo pasó entre el impulso y la acción, aunque haya sido «nada».}

\begin{infoband}{milcVerde}


\textbf{En tu cuaderno (Actividad 1 · IDENTIFICA).} Título: «Actividad 1. Tres casillas, tres veces». \textbf{Formato:} tres filas con tres columnas (sentí · me dieron ganas de · hice) más el tiempo que pasó. \textbf{Extensión:} media página. \textbf{Sabes que terminaste cuando} en alguna fila las columnas 2 y 3 son \textbf{distintas} ---ahí ya usaste el espacio, aunque no te dieras cuenta---. Esta página es tuya: \textbf{nadie más tiene que leerla si tú no quieres}.
\end{infoband}

\puente{Ya viste tus casos. Ahora las cuatro claves: qué es cada cosa, qué le pasa a una emoción con el tiempo, por qué nombrar frena, y qué hace realmente un fogón.}

\milcestacion{milcTurquesa}{Fase 2 · Sistematización: entender el espacio}

\begin{softbox}{milcTurquesa}{milcGris}{Cuatro claves sobre el espacio entre sentir y hacer}
Cuatro claves para entender por qué ese espacio existe, por qué a veces se cierra del todo y qué se puede hacer para que no se cierre.
\end{softbox}

\milcpaso{1}{milcTurquesa}{\textbf{Emoción, impulso y conducta son tres cosas, y solo una se elige.} La \textbf{emoción} llega sin permiso: no se discute ni se prohíbe, y no hay emociones malas. El \textbf{impulso} también llega solo: es lo que la emoción \emph{propone} ---la rabia propone golpear o gritar, el miedo propone huir o esconderse, la vergüenza propone desaparecer---. Y la \textbf{conducta} es lo que finalmente haces. \textbf{Solo la tercera se elige}, y por eso solo de la tercera te pueden pedir cuentas. \emph{Decir «me hizo enojar y por eso le pegué» junta las tres para que la última parezca inevitable. No lo es.}}
\milcpaso{2}{milcTurquesa}{\textbf{Una emoción sube, llega a un tope y baja ---si no la alimentas.} Ninguna rabia se queda en el máximo para siempre: sube rápido, se sostiene un rato y va bajando sola. Lo que la mantiene arriba es lo que \emph{haces} con ella: repetirte la escena, buscar a alguien que te dé la razón, releer el mensaje veinte veces, quedarte mirando la conversación. \textbf{A eso se le llama rumiar, y es echarle leña.} La consecuencia práctica es enorme: si aguantas sin actuar mientras la ola baja, \textbf{decides desde otro lugar}, y casi siempre decides mejor.}
\milcpaso{3}{milcTurquesa}{\textbf{Nombrar es frenar, y eso se puede medir.} No es un consejo bonito: cuando alguien le pone palabras a lo que siente, \textbf{baja la actividad de la amígdala} ---la que dispara la alarma--- y \textbf{sube la de la corteza prefrontal}, la parte que decide (Lieberman y otros, 2007). Por eso el muestrario del momento anterior es una herramienta y no una lista: decir por dentro \emph{«esto es rabia, y va como en 8»} hace en tu cabeza algo distinto que decir \emph{«estoy que exploto»}. \emph{Ponerle nombre no apaga la emoción: le baja el volumen lo suficiente para que quepa una decisión.}}
\milcpaso{4}{milcTurquesa}{\textbf{Un fogón es un espacio, no un sermón.} Lo que hace el nak chak misak no es dar la respuesta: es \textbf{obligar a un rodeo}. Lo que pasó en el camino se cuenta junto al fuego, se escucha a los mayores, y solo entonces se decide. Tu fogón funciona igual: no es un lugar donde te calmes por arte de magia, es \textbf{un rodeo obligatorio} entre el impulso y la acción. Puede ser un lugar (el patio, la escalera, la ventana), un gesto (respirar contando, tomar agua, salir a caminar) o una persona. \emph{Lo único que no puede ser es «pensar antes de actuar»: eso no es un fogón, es un deseo.}}

\begin{drawbox}{milcTurquesa}{Qué pasa en esos segundos, paso a paso}
\begin{pasoslista}
\item \textbf{Pasa algo.} No lo elegiste: un comentario, una nota, un mensaje, una mirada.
\item \textbf{Llega la emoción.} Tampoco la elegiste, y llegó por algo. El cuerpo avisa primero: pecho, cara, manos.
\item \textbf{Aparece el impulso.} La emoción propone una salida rápida. \emph{Proponer no es ordenar.}
\item \textbf{Aquí está el espacio.} Nombras lo que sientes, haces tu pausa, dices tu frase. Dura segundos, y es todo lo que hace falta.
\item \textbf{Eliges la conducta.} La misma emoción, otra salida. \emph{Esta parte sí lleva tu firma.}
\end{pasoslista}
\end{drawbox}

\begin{softbox}{milcMostaza}{milcGris}{Errores comunes y su corrección}
\textbf{«Es que soy así»:} el temperamento explica la velocidad de tu ola, no lo que haces con ella. \textbf{Confundir aguantarse con regular:} tragarse todo no ensancha el espacio, lo aplaza ---y suele estallar con quien menos tiene que ver---. \textbf{Creer que la pausa es dejarse:} no responder \emph{ya} no es perder; es decidir cuándo. \textbf{Pausas imposibles:} «respiro veinte minutos» no sirve en la mitad de una clase; que quepa en el lugar real. \textbf{Rumiar creyendo que estás pensando:} repetirte la escena no es analizarla, es echarle leña.
\end{softbox}

\begin{infoband}{milcTurquesa}


\textbf{En tu cuaderno (Actividad 2 · EXPLICA).} Título: «Actividad 2. El espacio entre sentir y hacer». \textbf{Formato:} los cinco pasos del recuadro aplicados a \textbf{una} de tus tres situaciones de la Actividad 1, y debajo dos líneas sobre qué le echó leña a esa emoción. \textbf{Extensión:} media página. \textbf{Sabes que terminaste cuando} puedes decir en voz alta, sobre tu propio caso, \textbf{cuál fue el impulso y cuál fue la decisión} ---y que no eran lo mismo---.
\end{infoband}

\puente{Ya entiendes el mecanismo. Es hora de construir el tuyo, porque en la mitad de la rabia nadie se acuerda de la teoría: se acuerda de un procedimiento ensayado.}

\milcestacion{milcMagenta}{Fase 3 · Praxis: construir tu fogón}

\begin{softbox}{milcMagenta}{milcGris}{Cómo se arma un procedimiento que funcione cuando estás alterado}
En la mitad de la rabia nadie se acuerda de lo que aprendió: se acuerda de lo que \textbf{ensayó}. Por eso el producto de hoy es un procedimiento corto, tuyo, escrito y practicado antes de necesitarlo. Tres pasos y un ensayo.
\end{softbox}

\milcpaso{1}{milcMagenta}{\textbf{La señal (qué te avisa a ti).} Cada persona tiene un aviso propio y físico, y casi siempre llega \textbf{antes} que el pensamiento: la mandíbula apretada, el calor en la cara, las manos frías, la voz que sube, el estómago cerrado, las ganas de llorar. Escribe \textbf{la tuya}, en concreto. \emph{Si no reconoces la señal, el procedimiento arranca tarde ---y tarde es después del portazo---.}}
\milcpaso{2}{milcMagenta}{\textbf{La pausa (qué haces exactamente).} Tiene que ser \textbf{corta, posible y del tamaño del lugar}: respirar contando hasta cuatro y soltar en seis, tomar agua, pedir permiso para salir un minuto, apretar las manos bajo el pupitre, dejar el celular boca abajo y no responder. \textbf{Escribe una para el salón y otra para la casa}, porque no son el mismo lugar ni tienen las mismas reglas.}
\milcpaso{3}{milcMagenta}{\textbf{La frase (qué te dices, por dentro).} Corta, tuya y creíble. \emph{«Esto es rabia y va en 8; en diez minutos va en 4»}. \emph{«No tengo que responder ya»}. \emph{«Si contesto ahora, mañana me toca arreglarlo»}. Evita las que no te crees: «no pasa nada» es mentira cuando sí pasa, y tu cabeza lo sabe.}
\milcpaso{4}{milcMagenta}{\textbf{El ensayo (en frío, con un compañero).} Uno describe una situación cotidiana que le suele dar rabia ---sin contar nada privado--- y el otro practica los tres pasos \textbf{en voz alta}: nombra la emoción, dice qué pausa haría, dice su frase. Luego cambian. \emph{Ensayarlo en frío es lo que hace que aparezca en caliente:} el procedimiento que no se dijo nunca en voz alta no llega cuando hace falta.}

\begin{drawbox}{milcMagenta}{Antes de cerrar tu fogón}
\begin{checklista}
\item Tu señal está escrita en concreto y es \textbf{física}, no un pensamiento («me dan ganas de gritar» no es la señal; «la mandíbula apretada» sí).
\item Tienes \textbf{dos pausas}: una que cabe en el salón y otra para la casa.
\item Tu frase es corta y te la crees. La leíste en voz alta y no te sonó falsa.
\item Ensayaste los tres pasos con un compañero, en frío y en voz alta.
\item Sabes cuál es tu «leña»: eso que haces y que mantiene la emoción arriba (releer, contarlo cinco veces, buscar quién te dé la razón).
\item Escribiste qué vas a hacer cuando el procedimiento \textbf{falle}, porque va a fallar alguna vez. Fallar no lo invalida: se vuelve a usar.
\end{checklista}
\end{drawbox}

\begin{softbox}{milcMagenta}{milcGris}{Plantilla de trabajo}
\textbf{Plantilla del fogón (cabe en una tarjeta).} \emph{Mi señal:} \_\_\_\_ · \emph{Mi pausa en el salón:} \_\_\_\_ · \emph{Mi pausa en la casa:} \_\_\_\_ · \emph{Mi frase:} «\_\_\_\_» · \emph{Mi leña (lo que no debo hacer):} \_\_\_\_ \\[1mm]
\textbf{Ejemplo.} \emph{Señal:} «se me calienta la cara y hablo más duro». \emph{Pausa en el salón:} «bajo la mirada al cuaderno y respiro cuatro veces contando». \emph{Pausa en la casa:} «salgo al patio dos minutos». \emph{Frase:} «no tengo que responder ya». \emph{Leña:} «releer el mensaje una y otra vez».
\end{softbox}

\begin{infoband}{milcMagenta}


\textbf{En tu cuaderno (Actividad 3 · CREA).} Título: «Actividad 3. Mi fogón». \textbf{Formato:} la tarjeta de cinco campos + una línea sobre cómo te fue en el ensayo. \textbf{Extensión:} media página, y cópiala también en una hoja pequeña que puedas cargar. \textbf{Sabes que terminaste cuando} podrías ejecutar los tres pasos \textbf{sin leer la tarjeta}, porque ya los dijiste en voz alta.
\end{infoband}

\puente{El producto es ese procedimiento. \emph{Un plan que no se ensayó no existe: existe la intención de tenerlo.}}

\milcestacion{milcMostaza}{Producto de la sesión}
\textbf{Mi fogón: señal, pausa y frase, ensayadas antes de necesitarlas}.
\begin{softbox}{milcMostaza}{milcGris}{Criterios de calidad}
(1) La señal es física y concreta, y la reconocerías en el momento en que aparece. (2) Hay dos pausas distintas, una posible en el salón y otra en la casa, ambas cortas. (3) La frase es propia, corta y creíble; no es una consigna copiada del tablero. (4) El procedimiento se ensayó en voz alta con un compañero, en frío. (5) Está identificada la «leña» propia y escrito qué hacer cuando el procedimiento falle.
\end{softbox}

\puente{Antes de cerrar, revisa si tu pausa es de verdad posible en el salón, en la casa y en la calle. Si solo sirve en un lugar, todavía no está listo.}

\milcestacion{milcVino}{Evaluación liberadora · cinco dimensiones}

\begin{softbox}{milcVino}{milcGris}{Sentido de la evaluación}
Evaluar no es solo revisar una nota. Es reconocer qué aprendí, qué cambió en mí, qué emociones manejé, cómo me relaciono con la ciudad y la comunidad, y qué le dejaré a quien viene detrás.
\end{softbox}

\textbf{Mi reflexión liberadora en cinco dimensiones.} Completa cada frase iniciada:

\small
\begin{tabularx}{\textwidth}{>{\bfseries\RaggedRight\arraybackslash}p{3.4cm}Y>{\RaggedRight\arraybackslash}p{4.6cm}}
\rowcolor{milcVino}\color{white}Dimensión & \color{white}\bfseries Mi respuesta (frame starter) & \color{white}\bfseries Algo que descubrí\\
1. Personal & Lo que más cambió en mí fue \linead{6.5cm} & \linead{4.2cm}\\
2. Emocional & La emoción más fuerte fue \linead{2.5cm} y la regulé \linead{3cm} & \linead{4.2cm}\\
3. Ciudadana & En democracia esto importa porque \linead{6cm} & \linead{4.2cm}\\
4. Local (Cartago) & En mi barrio sería útil para \linead{6.5cm} & \linead{4.2cm}\\
5. Intergeneracional & A mi abuela/abuelo le diría \linead{6.5cm} & \linead{4.2cm}\\
\end{tabularx}
\normalsize

\begin{infoband}{milcVino}
\textbf{Para la próxima sustentación.} Vuelve a esta tabla en seis meses. Verás que las respuestas cambian — no porque lo aprendido cambie, sino porque tú cambias. La evaluación liberadora es una conversación contigo a través del tiempo.
\end{infoband}


\puente{Cerramos con tres voces: lo que el impulso no deja ver, por qué la demora es el mejor remedio de la ira, y qué le hace al mundo compartido lo que decides en ese segundo.}

\milcestacion{milcGradeDark}{Triángulo de pensamiento · cierre filosófico}

\begin{softbox}{milcVino}{milcGris}{Dussel · lente del nosotros}
\textit{«Analéctico quiere indicar el hecho real humano por el que todo hombre, todo grupo o pueblo se sitúa siempre más allá (aná-) del horizonte de la totalidad.»}

{\footnotesize\itshape\color{milcNegro!70}--- Enrique Dussel · Filosofía de la liberación (1977), §2.4}

El impulso tiene una trampa: hace que la situación parezca completa, cerrada, sin más salidas que la que él propone. «O respondo, o me dejo». \textbf{Ese horizonte cerrado es lo que Dussel llama totalidad}, y la pausa sirve exactamente para asomarse más allá de él: casi siempre hay una tercera salida que la rabia no dejaba ver, y con frecuencia aparece cuando entra en cuenta alguien más ---qué le pasa al otro, qué no sé de él---. \emph{El fogón misak hace eso: mete otras voces antes de la decisión.}

\textbf{Cuando siento que solo hay una salida, ¿a quién podría escuchar antes de decidir que efectivamente no hay otra?}
\end{softbox}
\linead{\linewidth}\\[1mm]
\linead{\linewidth}

\begin{softbox}{milcMostaza}{milcGris}{Estoicismo · Séneca · Sobre la ira, libro II (c. 45 d.C.) · cuidado interior}
\textit{«El mayor remedio para la ira es la demora.»}

{\footnotesize\itshape\color{milcNegro!70}--- Séneca · Sobre la ira, libro II (c. 45 d.C.)}

Séneca no dice que la ira esté mal ni que haya que tragársela. Dice algo mucho más práctico, y hace dos mil años: \textbf{lo que la desarma es el tiempo}. Tu pausa es exactamente esa demora, en versión de once años y de cuatro respiraciones. No es cobardía ni es dejarse: es negarle a la ira el derecho a firmar por ti. \emph{Y encaja con lo que aprendiste hoy: la ola baja sola si no le echas leña.}

\textbf{¿Qué cosa de las que hice esta semana no habría hecho si hubiera esperado diez minutos?}
\end{softbox}
\linead{\linewidth}\\[1mm]
\linead{\linewidth}

\begin{softbox}{milcTurquesa}{milcGris}{Floridi · lente de la infoesfera}
\textit{«El florecimiento de las entidades informacionales y de toda la infoesfera debe promoverse preservando, cultivando y enriqueciendo sus propiedades.»}

{\footnotesize\itshape\color{milcNegro!70}--- Luciano Floridi · La ética de la información (2013; trad.)}

Lo que haces en ese segundo casi nunca se queda en ese segundo: el mensaje enviado con rabia queda escrito, se reenvía, se guarda en pantallazos y sigue existiendo cuando a ti ya se te pasó. \textbf{La emoción es tuya y dura minutos; lo que publicas con ella queda en un espacio que compartimos todos.} Por eso Floridi habla de \emph{cuidar} ese espacio: la pausa no solo te protege a ti de un lío, protege el lugar donde viven las conversaciones de tu curso.

\textbf{De lo que escribí este mes cuando estaba alterado, ¿qué sigue existiendo hoy en la pantalla de alguien?}
\end{softbox}
\linead{\linewidth}\\[1mm]
\linead{\linewidth}

\begin{infoband}{milcNegro}
\textbf{Nota para el docente.} Las tres son citas textuales y llevan su referencia debajo. Puedes remitir a la obra en clase.
\end{infoband}

\milcestacion{milcVino}{Mi autoevaluación · marca una casilla por fila}

{\small
\setlength{\extrarowheight}{2.2pt}
\begin{tabularx}{\textwidth}{>{\RaggedRight\arraybackslash\bfseries}p{3.0cm}YYY}
\rowcolor{milcVino}\color{white}\bfseries Criterio & \color{white}\bfseries Lo logré & \color{white}\bfseries En proceso & \color{white}\bfseries Necesito apoyo\\[.6mm]
Comprensión del tema & \milcopcion{Explico las ideas con mis palabras.} & \milcopcion{Reconozco las ideas, falta precisión.} & \milcopcion{Necesito repasar lo básico.}\\[.8mm]
\rowcolor{milcGris}Distinguir emoción y acción & \milcopcion{Separo con claridad lo que sentí, lo que me dieron ganas de hacer y lo que hice.} & \milcopcion{A veces confundo el impulso con la emoción.} & \milcopcion{Todavía digo «me hizo enojar y por eso lo hice».}\\[.8mm]
Aplicación y producto & \milcopcion{Mi producto cumple los criterios.} & \milcopcion{Está armado, con ajustes pendientes.} & \milcopcion{Aún está en construcción.}\\[.8mm]
\rowcolor{milcGris}Reflexión 5 dimensiones & \milcopcion{Respondí cada una con honestidad.} & \milcopcion{Respondí algunas con detalle.} & \milcopcion{Mis respuestas son muy generales.}\\[.8mm]
Triángulo de pensamiento & \milcopcion{Conecté las 3 voces con mi trabajo.} & \milcopcion{Conecté al menos una.} & \milcopcion{No las relaciono con mi proceso.}\\[.8mm]
\end{tabularx}}

\vspace{3mm}
\textbf{Hoy entendí que:}\\[1mm]
\linead{\linewidth}\\[1mm]
\linead{\linewidth}

\textbf{Mi compromiso para la próxima sesión es:}\\[1mm]
\linead{\linewidth}\\[1mm]
\linead{\linewidth}

\begin{softbox}{milcMostaza}{milcGris}{Cita de despedida · Marco Aurelio}
\textit{``El obstáculo en el camino se vuelve el camino. Lo que estorba al camino se vuelve camino.''}\\[1mm]
Cada error de hoy, bien observado, se convierte en pieza del camino que pisarás mañana.
\end{softbox}

\small
\begin{tabularx}{\textwidth}{>{\bfseries}p{3.6cm}Y}
\rowcolor{milcGradeDark!12}Nombre completo: & \linead{10cm}\\
Fecha: & \linead{4cm} \quad Curso/grupo: \linead{4cm}\\
Mi firma: & \linead{9cm}\\
\end{tabularx}
\normalsize

\begin{infoband}{milcGradeDark}
\textbf{Para la próxima sesión.} Dos cosas. \textbf{Primero, usa el fogón tres veces y anótalas} ---aunque te haya salido mal---: qué pasó, qué señal apareció, si alcanzaste a hacer la pausa y qué hiciste al final. \emph{Las veces que falla enseñan más que las que salen bien, porque muestran en qué punto exacto se te cierra el espacio.} \textbf{Segundo, pregunta en tu casa dónde queda el fogón de tu familia:} ¿en qué lugar se hablan las cosas difíciles? ¿La mesa, la cocina, el andén, el camino al colegio, el patio? Pregúntale a alguien mayor si en su casa de infancia había un lugar así y qué se decidía allí. \textbf{Trae:} tus tres registros y el lugar que te contaron. \emph{En el momento siguiente vas a trabajar la respiración, que es la pausa más portátil que existe: no necesita lugar, ni permiso, ni que nadie se dé cuenta.}
\end{infoband}

\end{document}
